\backhead{FURTHER EXERCISES}{khred}
\begin{exlist}
\item A log gives a storage organ dry matter of 7.2 Mg DM\,ha$^{-1}$. Find the yield at 14 per cent moisture.\exans{8.37 Mg\,ha$^{-1}$}
\item Close the balance: $P=540$, $I=120$, $T=290$, $E_s=150$, $E_o=45$, $D=160$, $Q=5$ mm, $\Delta S=-6$ mm.\exans{residual 16 mm: a term is missing}
\item A daily log of 20 years has how many rows? An hourly log?\exans{7305; 175\,320}
\item Convert 0.85 g\,cm$^{-3}$ to kg\,m$^{-3}$ and 12 kg N\,ha$^{-1}$ to g N\,m$^{-2}$.\exans{850 kg\,m$^{-3}$; 1.2 g\,m$^{-2}$}
\item The transpiration efficiency is 2.1 kg DM\,m$^{-3}$ and the transpiration 280 mm. Find the dry matter.\exans{5.88 Mg\,ha$^{-1}$}
\end{exlist}
\backhead{ESSAY QUESTIONS}{khred}
{\small\itshape Answer each question in about one page, with equations, a sketch and at least one number from this chapter.}\par\vspace{3pt}
\begin{exlist}
\item Describe the structure of the Daisy output and how it is read into SWAFL.
\item Discuss the checks that must be made before a log is used.
\item Explain why daily figures are drawn only from some experiments.
\item Discuss the conversion of dry matter into yield and water productivity.
\item Explain how a failed run is recognised and what can be learnt from it.
\item Describe how you would present the output of a Daisy experiment in a report.
\end{exlist}
\dykband{EXERCISES IN SWAFL}
{\small Open SWAFL (desktop or online) with the benchmark project. The chapter of the \href{https://github.com/halilunaziru73-creator/Coupled-Water-Fluxes-Pore-Scale-Dynamics-Precision-Irrigation-Mulching-Feedback-Loops/raw/main/Software/v1.1.0/SWAFL_Desktop/docs/manual_v1.1.0/SWAFL_Manual_v1.1.0.pdf}{SWAFL manual} that explains each task is given in brackets; record your values in a table and compare them with what is expected (\emph{Check}).}\par\vspace{3pt}
\begin{exlist}
\item {\color{khblue}[Manual, Chapter 30]} Open the raw output of a run and identify the header, names, units and rows.\par\noindent{\small\itshape Check: as in the figure of this chapter.}
\item {\color{khblue}[Manual, Chapter 30]} Open the efficiency indices of an experiment and compute one by hand.\par\noindent{\small\itshape Check: the hand value agrees.}
\item {\color{khblue}[Manual, Chapter 30]} Find where the results are saved and open one log file in a text editor.\par\noindent{\small\itshape Check: the folder holds every log.}
\item {\color{khblue}[Manual, Chapter 30]} Open \emph{Why a run gave no results} and explain one failure.\par\noindent{\small\itshape Check: the line of the set-up is shown.}
\end{exlist}
\backhead{OBJECTIVE TYPE QUESTIONS}{khred}
\begin{exlist}
\item A Daisy log file has the extension\opt{.dai}{.dlf}{.csv}{.log}
\item The yield at 15 per cent moisture of 6 Mg DM\,ha$^{-1}$ is\opt{5.1}{6.0}{7.06}{6.9} Mg\,ha$^{-1}$
\item A residual of 20 mm in a yearly balance indicates\opt{a good closure}{a missing term}{a wet year}{a dry year}
\item A daily log over 20 years has about\opt{365}{7300}{20\,000}{175\,000} rows
\item Transpiration efficiency is expressed in\opt{kg\,m$^{-3}$}{mm}{kg\,ha$^{-1}$}{dimensionless}
\end{exlist}
\objkey{1.\ (b),\quad 2.\ (c),\quad 3.\ (b),\quad 4.\ (b),\quad 5.\ (a)}
